\subsection{CHARTS AND POINTS}

\subsubsection{TURN RECORD CHART}

The TURN RECORD CHART is used for time-keeping. The Turn Marker Counter is placed in the correctly numbered square on the chart to indicate the current Turn. As the game progresses, and each Turn is finished, the Confederate Player advances the marker on the track into the next-numbered box.

Also, any new units that are to be introduced onto the Mapboard after the game begins should be placed in the correctly-numbered box to indicate the Turn of their arrival on the Mapboard. Their Turn of arrival and place of arrival is given in the individual scenarios.

\subsubsection{POINT RECORD}

Both sides use the Point Record, the Union Player indicating his current point levels through the use of the three blue Point counters, and the Confederate Player indicating his current point levels through use the three gray point counters.

\textbf{Procedure:} The three tracks of the POINT RECORD are number for ones, tens and hundreds, and the three Point counters are placed in appropriate boxes to indicate the player's Victory Point count.

\textit{For example, a score of 123 would be indicated by placing one counter in the "100" box, one counter in the "20" box and one counter in the "3" box.}

\subsubsection{HOW TO WIN}

Players keep track on the POINT RECORD of the Victory Points that they accumlate during the course of the game. These are compared at game's end when all Turns are complete (different scenarios last varying numbers of Turns). The Level of Victory is defined for each scenario, but basically the winner is that side which has accumlated the most points.

Point values are listed on the VICTORY POINTS TABLE on the back page of the Rules Folder. These point values are used for all scenarios.

\subsubsection{AN EXPLANATION OF VICTORY POINTS}

Players gain Victory Points for their side through a variety of methods, explained below.

\begin{enumerate}
  \item \textbf{CASUALTIES:}

  Victory Points can be gained by eliminating the other side's Combat Factors. For instance, the loss of one Union Combat Factor of Infantry gives the Confederate player 4 Victory Points. This also applies to the destruction and/or capture of enemy Supply or Wagon counters (a player who destroys his own Supply or Wagon counters denies their use to the enemy, but the enemy still gets the Victory Points for their Destruction). Note that the Union, having the greater overall resources, can affor to take greater losses, as each individual is worth less Victory Points to the enemy

  \item \textbf{GEOGRAPHICAL:}

  Victory Points can be gained by capturing key points from the enemy, or by denying him the use of his railroads. Note that a rail line can be cut at any point, but must be still cut at the end of a Turn to count. A rail line is cut if an enemy occupies a hex of the line, or, if it is destroyed at some point. The capture of the towns of Romney, Franklin, and Moorefield are worth points to the Confederate if they are entered at any point during the game, but they are only worth Victory Points once per game (they are also only worth points in scenarios where they contain a Union garrison at the start of the game). Staunton, Charlottesville, and Harper's Ferry are worth points only if held at the end of the game. Also not that these three towns are worth a varying number of points, depending on the type and supplied status of the units in them.

  \item \textbf{VICTORIES:}

  The Shenandoah region was very close to the news centers for both the Union and Confederacy. Therefor, a very small victory or defeat was blown up to headline proportions, and had considerable propaganda value. Therefore, players can get Victory Points for forcing an enemy to retreat (more points if he outnumbers you), or by totally wiping out some enemy force (the capture of their colors, guns, etc, gives the victor suitable trophies of victory).

  \item \textbf{COMMMITMENT OF RESOURCES:}

  The North, having greater resources in men and materials could afford to commit more to the Shenandoah than the South could. Since the Shenandoah was basically a secondary theater of operations, both sides were constantly trying to hold it with a minimum force while, at the same time, trying to tie up as large a proportion of the enemy's resources as possible. Thus, it is possible for the Union to outnumber the Confederates, and yet still have a smaller committment of resources. The entrance of new Supply and/or Wagon units onto the Mapboard represents the commitment of more resources, and thus gives the opposition Victory Points. The overall size of the Combat forces committed also reflects the resources used by both sides. Each scenario gives a Resource Factor for each side. At the end of each Turn, the smaller Resource Factor is subtracted from the larger, and the difference is given to the side with the smaller Resource Factor in the form of Victory Points.

  \textit{For example, if the Confederates had 5 Resource Factors to the Union's 7 Resource Factors, the Confederates would get two Victory Points.}

  The entry of new units onto the Mapboard will increase a side's Resource Factor, as given in each scenario. A side can decrease it's resource factor by "one" for each 50 Victory Points worth of Combat Factors it exits from the Mapboard at a designated exit hex.

  \textit{For example, if the Union Player exited 11 Infantry Combat Factors and one Cavalry Combat Factor (11 x 4 = 44 + 6 = 50 Victory Points), he would reduce his Resource Factor by "one".}

  Supply and Wagon counters cannot be exited for Victory Points, unless stated in the scenario. Partisan units cannot be exited for reducing the Confederate Resource Factor. General Officer counters \underline{can} be exited to reduce the Resource Factor. Resource Factors are not reduced by casualties.

  \item \textbf{THE VALLEY:}

  The Shenandoah region itself was a resource to the Confederacy. The "Garden of Virginia" played a major role in feeding Lee's Army of Northern Virginia. This was finally realized by the Union commanders rather late in the war (it can only be used during scenarios \#10 - \#16), and they initiated a policy of systematic devastation. Each seven-hexes Devastated Area counts for two Victory Points for the Union.

  \item \textbf{SPECIAL CONDITIONS:}

  Some of the scenarios have special conditions for acquiring Victory Points that apply only to that scenario. These are listed in the scenario.
\end{enumerate}

{\large \textbf{END OF BASIC GAME RULES}}
